% ------------------------------------------------------------ Ký hiệu và chữ viết tắt
\chapter*{DANH MỤC KÝ HIỆU VÀ CHỮ VIẾT TẮT}
\addcontentsline{toc}{chapter}{Danh mục ký hiệu và chữ viết tắt}

\begingroup
\renewcommand{\arraystretch}{1.18}
\noindent\textbf{Ký hiệu}

\smallskip
\noindent\begin{tabularx}{\linewidth}{@{}p{3.1cm}X@{}}
$\R^n$ & không gian Euclid $n$ chiều; $n$ là số chiều của dữ liệu đầu vào\\
$x(i)$ & thành phần thứ $i$ của vectơ $x\in\R^n$\\
$x_k,\ y_k$ & điểm dữ liệu thứ $k$ và nhãn của nó, $y_k\in\{-1,+1\}$, $k=1,\dots,N$\\
$N$ & số điểm dữ liệu huấn luyện\\
$(t)_+$ & hàm bản lề $\max\{0,t\}$\\
$\intr C,\ \conv C$ & phần trong và bao lồi của tập $C$\\
$\Pi_C(x)$ & hình chiếu của $x$ lên tập lồi đóng $C$\\
$\Omega,\ \Omega_k$ & miền xác định và các vùng (đa diện) của một phân hoạch\\
$\B$ & tập (biên) tuyến tính từng phần\\
$\varphi(x)$ & ánh xạ đặc trưng, $\varphi:\R^n\to\R^M$; $\varphi_m$ là đặc trưng thứ $m$\\
$M$ & số đặc trưng; $M/n$ là số đặc trưng trên mỗi chiều\\
$D$ & số đặc trưng bản lề, $D=M-n$\\
$p_m,\ q_m$ & tham số của siêu phẳng bản lề $p_m^\T x+q_m=0$ của đặc trưng thứ $m$\\
$w,\ w_0$ & vectơ trọng số và hệ số chặn của bộ phân loại\\
$\gamma$ & hằng số chính quy hóa (trọng số của hàm mất mát) trong SVM\\
$\kappa(x,z)$ & hạt nhân, $\kappa(x,z)=\varphi(x)^\T\varphi(z)$\\
$\alpha$ & vectơ biến đối ngẫu (nhân tử Lagrange)\\
$e_k$ & biến bù của điểm dữ liệu thứ $k$\\
$N_s$ & số vectơ hỗ trợ\\
$h_i$ & hàm đóng góp của biến thứ $i$ trong mô hình cộng tính\\
$\mathbf 1$ & vectơ có mọi thành phần bằng $1$\\
$\sgn(t)$ & hàm dấu\\
\end{tabularx}

\bigskip
\noindent\textbf{Chữ viết tắt}

\smallskip
\noindent\begin{tabularx}{\linewidth}{@{}p{3.1cm}X@{}}
AUC & diện tích dưới đường cong ROC (area under the curve)\\
C-SVM & SVM lề mềm với hàm mất mát bản lề\\
ĐCX-CB & độ chính xác cân bằng (balanced accuracy)\\
F1 & trung bình điều hòa của độ chuẩn xác (precision) và độ nhạy (recall)\\
GHH & siêu phẳng bản lề tổng quát (generalized hinging hyperplane)\\
HH & siêu phẳng bản lề (hinging hyperplane)\\
Ik-SVM & SVM với hạt nhân giao (intersection kernel)\\
KKT & điều kiện Karush--Kuhn--Tucker\\
L-BFGS & phương pháp tựa Newton BFGS với bộ nhớ hạn chế\\
kNN & phương pháp $k$ láng giềng gần nhất\\
LS-SVM & SVM bình phương tối thiểu (least squares SVM)\\
MLP & mạng nơ-ron truyền thẳng nhiều lớp (multilayer perceptron)\\
PWL & tuyến tính từng phần (piecewise linear)\\
PWL-SVM & SVM với ánh xạ đặc trưng tuyến tính từng phần\\
RBF & hàm cơ sở bán kính; hạt nhân Gauss\\
ReLU & đơn vị tuyến tính chỉnh lưu, $\mathrm{ReLU}(t)=\max\{0,t\}$\\
ROC & đường đặc trưng hoạt động (receiver operating characteristic)\\
SVM & máy vectơ hỗ trợ (support vector machine)\\
\end{tabularx}
\endgroup
